% ============================================================================
\section{Governance and standards}
% ============================================================================

\begin{frame}{OWASP Top 10 for LLM Applications}
  \begin{itemize}\itemsep0.5em
    \item[\textcolor{KAred}{$\blacktriangleright$}] A community-maintained list of the most critical risks in LLM applications, widely used as a checklist in industry.
    \item[\textcolor{KAred}{$\blacktriangleright$}] Prompt injection is ranked first (LLM01). The threats in this lecture map onto the list:
      \begin{itemize}\footnotesize\itemsep0.25em
        \item LLM01 Prompt Injection and LLM02 Sensitive Information Disclosure.
        \item LLM06 Excessive Agency and LLM04 Data and Model Poisoning.
        \item LLM05 Improper Output Handling and LLM08 Vector and Embedding Weaknesses.
      \end{itemize}
    \item[\textcolor{KAred}{$\blacktriangleright$}] The list gives teams a shared vocabulary and a baseline of risks to test against before shipping.
  \end{itemize}
  \imgsrc{OWASP Top 10 for LLM Applications (2025), \href{https://genai.owasp.org/}{genai.owasp.org}.}
\end{frame}

\begin{frame}{NIST AI Risk Management Framework}
  \begin{itemize}\itemsep0.5em
    \item[\textcolor{KAred}{$\blacktriangleright$}] A voluntary US framework for managing AI risk across the lifecycle, organized around four functions:
      \begin{itemize}\itemsep0.25em
        \item \textbf{Govern:} policies, roles, and accountability for AI risk.
        \item \textbf{Map:} identify the context and where harm could arise.
        \item \textbf{Measure:} assess and track risks with tests and metrics.
        \item \textbf{Manage:} prioritize, mitigate, and monitor over time.
      \end{itemize}
    \item[\textcolor{KAred}{$\blacktriangleright$}] A companion Generative AI Profile (2024) adapts these functions to generative and agentic systems.
    \item[\textcolor{KAred}{$\blacktriangleright$}] It addresses how an organization governs the risk, complementing the specific vulnerabilities OWASP enumerates.
  \end{itemize}
\end{frame}

\begin{frame}{Provider Policies and Deployment Gates}
  \begin{itemize}\itemsep0.5em
    \item[\textcolor{KAred}{$\blacktriangleright$}] Model providers publish safety policies and tie the release of high-capability, high-autonomy features to safety evaluations passing first.
    \item[\textcolor{KAred}{$\blacktriangleright$}] \textbf{Capability thresholds:} the more autonomous or consequential the deployment, the stronger the required safeguards and review before launch.
    \item[\textcolor{KAred}{$\blacktriangleright$}] \textbf{Staged rollout:} start with a limited release, monitor real usage, then widen, rather than shipping full autonomy to everyone at once.
    \item[\textcolor{KAred}{$\blacktriangleright$}] A deployment gate checks for the same technical safeguards covered in this lecture.
  \end{itemize}
\end{frame}
